\documentclass[10pt,a4paper]{amsart}
\usepackage{amsmath,amssymb,amsthm,graphicx}
% Retained minimal preamble for the Stacks Project chapter more-algebra.tex (More on Algebra).
% The chapter file pinned-original-source.tex carries only \input{preamble} on its first line, so no
% preamble text lives inside the chapter itself.  The Stacks Project uses one single shared preamble for
% every chapter; the snapshot stages it identically as sources/stacks-chow/preamble.tex and
% sources/stacks-etale-cohomology/preamble.tex.  The block below is copied mechanically, line for line,
% from sources/stacks-chow/preamble.tex lines 14-18 (the Xy-pic setup) and 154-217 (the native theorem
% declarations together with the mathematical macros).  Nothing was added, removed or redefined; the
% wrapper class and the amsmath/amssymb/amsthm/graphicx base packages are supplied by the assembler.
% Chapter: more-algebra.tex -- More on Algebra.
% For commutative diagrams we use Xy-pic
\usepackage[all]{xy}

% We use 2cell for 2-commutative diagrams.
\xyoption{2cell}

% Theorem environments.
%
\theoremstyle{plain}
\newtheorem{theorem}[subsection]{Theorem}
\newtheorem{proposition}[subsection]{Proposition}
\newtheorem{lemma}[subsection]{Lemma}

\theoremstyle{definition}
\newtheorem{definition}[subsection]{Definition}
\newtheorem{example}[subsection]{Example}
\newtheorem{exercise}[subsection]{Exercise}
\newtheorem{situation}[subsection]{Situation}

\theoremstyle{remark}
\newtheorem{remark}[subsection]{Remark}
\newtheorem{remarks}[subsection]{Remarks}

\numberwithin{equation}{subsection}

% Macros
%
\def\lim{\mathop{\mathrm{lim}}\nolimits}
\def\colim{\mathop{\mathrm{colim}}\nolimits}
\def\Spec{\mathop{\mathrm{Spec}}}
\def\Hom{\mathop{\mathrm{Hom}}\nolimits}
\def\Ext{\mathop{\mathrm{Ext}}\nolimits}
\def\SheafHom{\mathop{\mathcal{H}\!\mathit{om}}\nolimits}
\def\SheafExt{\mathop{\mathcal{E}\!\mathit{xt}}\nolimits}
\def\Sch{\mathit{Sch}}
\def\Mor{\mathop{\mathrm{Mor}}\nolimits}
\def\Ob{\mathop{\mathrm{Ob}}\nolimits}
\def\Sh{\mathop{\mathit{Sh}}\nolimits}
\def\NL{\mathop{N\!L}\nolimits}
\def\CH{\mathop{\mathrm{CH}}\nolimits}
\def\proetale{{pro\text{-}\acute{e}tale}}
\def\etale{{\acute{e}tale}}
\def\QCoh{\mathit{QCoh}}
\def\Ker{\mathop{\mathrm{Ker}}}
\def\Im{\mathop{\mathrm{Im}}}
\def\Coker{\mathop{\mathrm{Coker}}}
\def\Coim{\mathop{\mathrm{Coim}}}

% Boxtimes
%
\DeclareMathSymbol{\boxtimes}{\mathbin}{AMSa}{"02}

%
% Macros for moduli stacks/spaces
%
\def\QCohstack{\mathcal{QC}\!\mathit{oh}}
\def\Cohstack{\mathcal{C}\!\mathit{oh}}
\def\Spacesstack{\mathcal{S}\!\mathit{paces}}
\def\Quotfunctor{\mathrm{Quot}}
\def\Hilbfunctor{\mathrm{Hilb}}
\def\Curvesstack{\mathcal{C}\!\mathit{urves}}
\def\Polarizedstack{\mathcal{P}\!\mathit{olarized}}
\def\Complexesstack{\mathcal{C}\!\mathit{omplexes}}
% \Pic is the operator that assigns to X its picard group, usage \Pic(X)
% \Picardstack_{X/B} denotes the Picard stack of X over B
% \Picardfunctor_{X/B} denotes the Picard functor of X over B
\def\Pic{\mathop{\mathrm{Pic}}\nolimits}
\def\Picardstack{\mathcal{P}\!\mathit{ic}}
\def\Picardfunctor{\mathrm{Pic}}
\def\Deformationcategory{\mathcal{D}\!\mathit{ef}}

% --- appended 2026-10-12 (Chief): the project own definitions of the reference/slogan/history
% display environments, copied mechanically from the shared Stacks preamble lines 8-12
% (\usepackage{verbatim} plus the three \newenvironment lines, which the chapter preamble keeps
% verbatim from the project preamble). Nothing added, removed or redefined: these environments
% carry no mathematical content and the project itself typesets them as comments. Without these
% three lines a counted statement whose own text uses them fails with Environment slogan undefined.
\usepackage{verbatim}
\newenvironment{reference}{\comment}{\endcomment}
\newenvironment{slogan}{\comment}{\endcomment}
\newenvironment{history}{\comment}{\endcomment}
\usepackage{cleveref}
\crefname{originaldefinition}{original definition}{original definitions}
\Crefname{originaldefinition}{Original definition}{Original definitions}
\makeatletter
\newlabel{definition-perfect}{{0657}{source}{}{sourcepointer-1}{}}
\newlabel{definition-perfect@cref}{{[originaldefinition][0][]0657}{[1][0][]source}{}{}{}}
\makeatother
\crefname{originallemma}{original lemma}{original lemmas}
\Crefname{originallemma}{Original lemma}{Original lemmas}
\makeatletter
\newlabel{lemma-summands-perfect}{{066S}{source}{}{sourcepointer-2}{}}
\newlabel{lemma-summands-perfect@cref}{{[originallemma][0][]066S}{[1][0][]source}{}{}{}}
\makeatother
\crefname{originallemma}{original lemma}{original lemmas}
\Crefname{originallemma}{Original lemma}{Original lemmas}
\makeatletter
\newlabel{lemma-commutes-with-countable-sums}{{07LR}{source}{}{sourcepointer-3}{}}
\newlabel{lemma-commutes-with-countable-sums@cref}{{[originallemma][0][]07LR}{[1][0][]source}{}{}{}}
\makeatother
\crefname{originaldefinition}{original definition}{original definitions}
\Crefname{originaldefinition}{Original definition}{Original definitions}
\makeatletter
\newlabel{derived-definition-compact-object}{{07LS}{source}{}{sourcepointer-4}{}}
\newlabel{derived-definition-compact-object@cref}{{[originaldefinition][0][]07LS}{[1][0][]source}{}{}{}}
\makeatother
\crefname{originallemma}{original lemma}{original lemmas}
\Crefname{originallemma}{Original lemma}{Original lemmas}
\makeatletter
\newlabel{derived-lemma-morphisms-from-projective-complex}{{064B}{source}{}{sourcepointer-5}{}}
\newlabel{derived-lemma-morphisms-from-projective-complex@cref}{{[originallemma][0][]064B}{[1][0][]source}{}{}{}}
\makeatother
\crefname{originallemma}{original lemma}{original lemmas}
\Crefname{originallemma}{Original lemma}{Original lemmas}
\makeatletter
\newlabel{derived-lemma-subcategory-left-resolution}{{05T7}{source}{}{sourcepointer-6}{}}
\newlabel{derived-lemma-subcategory-left-resolution@cref}{{[originallemma][0][]05T7}{[1][0][]source}{}{}{}}
\makeatother
\crefname{originallemma}{original lemma}{original lemmas}
\Crefname{originallemma}{Original lemma}{Original lemmas}
\makeatletter
\newlabel{derived-lemma-projectors-have-images-triangulated}{{05QW}{source}{}{sourcepointer-7}{}}
\newlabel{derived-lemma-projectors-have-images-triangulated@cref}{{[originallemma][0][]05QW}{[1][0][]source}{}{}{}}
\makeatother
\crefname{originallemma}{original lemma}{original lemmas}
\Crefname{originallemma}{Original lemma}{Original lemmas}
\makeatletter
\newlabel{derived-lemma-representable-homological}{{0149}{source}{}{sourcepointer-8}{}}
\newlabel{derived-lemma-representable-homological@cref}{{[originallemma][0][]0149}{[1][0][]source}{}{}{}}
\makeatother
\crefname{originallemma}{original lemma}{original lemmas}
\Crefname{originallemma}{Original lemma}{Original lemmas}
\makeatletter
\newlabel{derived-lemma-derived-canonical-delta-functor}{{0152}{source}{}{sourcepointer-9}{}}
\newlabel{derived-lemma-derived-canonical-delta-functor@cref}{{[originallemma][0][]0152}{[1][0][]source}{}{}{}}
\makeatother
\crefname{originalsection}{original section}{original sections}
\Crefname{originalsection}{Original section}{Original sections}
\makeatletter
\newlabel{homology-section-truncations}{{0118}{source}{}{sourcepointer-10}{}}
\newlabel{homology-section-truncations@cref}{{[originalsection][0][]0118}{[1][0][]source}{}{}{}}
\makeatother
\title{A complex is perfect exactly when it is compact as an object of the derived category}
\author{The Stacks Project Authors}
\date{stacks-project:more-algebra.tex:local-snapshot-20261012; GNU FDL 1.2 or later}
\begin{document}
\maketitle
\setcounter{subsection}{1}
\noindent\textit{Editorial note.} All delivered mathematical text is the author's own, copied from the pinned Stacks Project chapter file more-algebra.tex; no statement or proof step was rewritten, reordered or removed. The counted span is the author environment and its own proof environment. References inside the span are handed to the assembler as navigation stubs that carry the canonical tag of the cited tag only; whether each of them is load-bearing is recorded on the stub itself.\par\medskip
\noindent\textit{Allowed standard prerequisites (editorial).}\begin{itemize}
\item The Hom functor out of a bounded above complex of projective modules is computed on the level of complexes, so its cohomology is the cohomology of a product complex. used at proof lines 20400-20417. Identifies the Hom groups out of the bounded complex of finite projective modules representing a perfect object with the cohomology of a product of Hom modules.
\item A compact object of the derived category has a map to it from a bounded complex of finite free modules which admits a one-sided inverse. invoked at proof lines 20470-20500. Produces the finite free complex L, the maps making K a summand of L, and the idempotent endomorphism of L.
\end{itemize}\medskip
\ifdefined\hypertarget\hypertarget{sourcepointer-1}{}\fi
\noindent\textit{Source reference.} Original definition 0657: sources/stacks-more-algebra/pinned-original-source.tex, source lines 18871--18871. The sentence at proof line 20399 recalls that perfect means quasi-isomorphic to a bounded complex of finite projective modules; the whole use of that representation, the product complex and the factoring argument stay inside the counted span.\par\medskip
\ifdefined\hypertarget\hypertarget{sourcepointer-2}{}\fi
\noindent\textit{Source reference.} Original lemma 066S: sources/stacks-more-algebra/pinned-original-source.tex, source lines 18963--18963. The conclusion at proof line 20506 is that K is perfect because it is a summand of the finite free complex L; the identification of K with the kernel of the complementary idempotent, which is what makes the summand statement applicable, is carried out inside the counted span.\par\medskip
\ifdefined\hypertarget\hypertarget{sourcepointer-3}{}\fi
\noindent\textit{Source reference.} Original lemma 07LR: sources/stacks-more-algebra/pinned-original-source.tex, source lines 20332--20332. The reference is the statement that a map out of a bounded complex of finite free modules into a direct sum factors through finitely many terms; this is exactly the compactness conclusion being proved, and its verification by the product complex argument is inside the counted span.\par\medskip
\ifdefined\hypertarget\hypertarget{sourcepointer-4}{}\fi
\noindent\textit{Source reference.} Original definition 07LS: sources/stacks-chow/tags.txt, source lines 1--1. The reference names the definition of compactness that the backward direction assumes at proof line 20441; that hypothesis and every deduction drawn from it stay inside the counted span.\par\medskip
\ifdefined\hypertarget\hypertarget{sourcepointer-5}{}\fi
\noindent\textit{Source reference.} Original lemma 064B: sources/stacks-chow/tags.txt, source lines 1--1. Used at proof line 20417 to replace Hom in the derived category out of a bounded above projective complex by the Hom complex; the subsequent computation of the product complex and the factoring of a map stay inside the counted span.\par\medskip
\ifdefined\hypertarget\hypertarget{sourcepointer-6}{}\fi
\noindent\textit{Source reference.} Original lemma 05T7: sources/stacks-chow/tags.txt, source lines 1--1. Used at proof lines 20441-20470 to produce the bounded complex L of finite free modules and the one-sided inverse; the splitting of L and the identification of K as its summand are inside the counted span.\par\medskip
\ifdefined\hypertarget\hypertarget{sourcepointer-7}{}\fi
\noindent\textit{Source reference.} Original lemma 05QW: sources/stacks-chow/tags.txt, source lines 1--1. Used at proof line 20490 to split L as the kernel of the idempotent together with the kernel of its complement; constructing the idempotent and identifying the relevant summand with K is done inside the counted span.\par\medskip
\ifdefined\hypertarget\hypertarget{sourcepointer-8}{}\fi
\noindent\textit{Source reference.} Original lemma 0149: sources/stacks-chow/tags.txt, source lines 1--1. A navigation pointer inside the forward direction; the argument that a map out of the perfect object factors through a finite direct sum is carried out inside the counted span.\par\medskip
\ifdefined\hypertarget\hypertarget{sourcepointer-9}{}\fi
\noindent\textit{Source reference.} Original lemma 0152: sources/stacks-chow/tags.txt, source lines 1--1. A navigation pointer for the manipulation of direct sums in the derived category; the termwise description of the direct sum used in the proof is inside the counted span.\par\medskip
\ifdefined\hypertarget\hypertarget{sourcepointer-10}{}\fi
\noindent\textit{Source reference.} Original section 0118: sources/stacks-chow/tags.txt, source lines 1--1. The proof refers to that section twice for the truncation conventions used when replacing a complex by a finite one; the truncations themselves and the estimates they feed are inside the counted span.\par\medskip
\begin{proposition}
\label{proposition-perfect-is-compact}
Let $R$ be a ring. For an object $K$ of $D(R)$ the following are equivalent
\begin{enumerate}
\item $K$ is perfect, and
\item $K$ is a compact object of $D(R)$.
\end{enumerate}
\end{proposition}

\begin{proof}
Assume $K$ is perfect, i.e., $K$ is quasi-isomorphic to a bounded complex
$P^\bullet$ of finite projective modules, see
Definition \ref{definition-perfect}. If $E_i$ is represented by the complex
$E_i^\bullet$, then $\bigoplus E_i$ is represented by the complex
whose degree $n$ term is $\bigoplus E_i^n$. On the other hand,
as $P^n$ is projective for all $n$ we have
$\Hom_{D(R)}(P^\bullet, K^\bullet) = \Hom_{K(R)}(P^\bullet, K^\bullet)$
for every complex of $R$-modules $K^\bullet$, see
Derived Categories,
Lemma \ref{derived-lemma-morphisms-from-projective-complex}.
Thus $\Hom_{D(R)}(P^\bullet, E^\bullet)$ is the cohomology of the complex
$$
\prod \Hom_R(P^n, E^{n - 1}) \to
\prod \Hom_R(P^n, E^n) \to
\prod \Hom_R(P^n, E^{n + 1}).
$$
Since $P^\bullet$ is bounded we see that we may replace the $\prod$
signs by $\bigoplus$ signs in the complex above. Since each $P^n$ is a finite
$R$-module we see that
$\Hom_R(P^n, \bigoplus_i E_i^m) = \bigoplus_i \Hom_R(P^n, E_i^m)$
for all $n, m$.
Combining these remarks we see that the map of
Derived Categories, Definition \ref{derived-definition-compact-object}
is a bijection.

\medskip\noindent
Conversely, assume $K$ is compact.
Represent $K$ by a complex $K^\bullet$ and consider the map
$$
K^\bullet
\longrightarrow
\bigoplus\nolimits_{n \geq 0} \tau_{\geq n} K^\bullet
$$
where we have used the canonical truncations, see
Homology, Section \ref{homology-section-truncations}.
This makes sense as in each degree the direct sum on the right is finite.
By assumption this map factors through a finite direct sum.
We conclude that $K \to \tau_{\geq n} K$ is zero for at least one $n$,
i.e., $K$ is in $D^{-}(R)$.

\medskip\noindent
Since $K \in D^{-}(R)$ and since every $R$-module is a quotient of a free
module, we may represent $K$ by a bounded above complex $K^\bullet$
of free $R$-modules, see
Derived Categories, Lemma \ref{derived-lemma-subcategory-left-resolution}.
Note that we have
$$
K^\bullet = \bigcup\nolimits_{n \leq 0} \sigma_{\geq n}K^\bullet
$$
where we have used the stupid truncations, see
Homology, Section \ref{homology-section-truncations}.
Hence by Lemma \ref{lemma-commutes-with-countable-sums} we see that
$1 : K^\bullet \to K^\bullet$ factors through
$\sigma_{\geq n}K^\bullet \to K^\bullet$ in $D(R)$.
Thus we see that $1 : K^\bullet \to K^\bullet$ factors as
$$
K^\bullet \xrightarrow{\varphi} L^\bullet \xrightarrow{\psi} K^\bullet
$$
in $D(R)$ for some complex $L^\bullet$ which is bounded and whose terms
are free $R$-modules. Say $L^i = 0$ for $i \not \in [a, b]$.
Fix $a, b$ from now on. Let $c$ be the largest integer $\leq b + 1$
such that we can find a factorization of $1_{K^\bullet}$ as above
with $L^i$ is finite free for $i < c$. We will show by induction that
$c = b + 1$. Namely, write $L^c = \bigoplus_{\lambda \in \Lambda} R$.
Since $L^{c - 1}$ is finite free we can find a finite subset
$\Lambda' \subset \Lambda$ such that $L^{c - 1} \to L^c$ factors
through $\bigoplus_{\lambda \in \Lambda'} R \subset L^c$. Consider the
map of complexes
$$
\pi :
L^\bullet
\longrightarrow
(\bigoplus\nolimits_{\lambda \in \Lambda \setminus \Lambda'} R)[-c]
$$
given by the projection onto the factors corresponding to
$\Lambda \setminus \Lambda'$ in degree $c$.
By our assumption on $K$ we see that, after possibly replacing $\Lambda'$ by
a larger finite subset, we may assume that $\pi \circ \varphi = 0$
in $D(R)$. Let $(L')^\bullet \subset L^\bullet$ be the kernel of $\pi$.
Since $\pi$ is surjective we get a short exact sequence of complexes,
which gives a distinguished triangle in $D(R)$ (see
Derived Categories, Lemma \ref{derived-lemma-derived-canonical-delta-functor}).
Since $\Hom_{D(R)}(K, -)$ is homological (see
Derived Categories, Lemma \ref{derived-lemma-representable-homological})
and $\pi \circ \varphi = 0$, we can find a morphism
$\varphi' : K^\bullet \to (L')^\bullet$ in $D(R)$ whose
composition with $(L')^\bullet \to L^\bullet$ gives $\varphi$.
Setting $\psi'$ equal to the composition of $\psi$ with
$(L')^\bullet \to L^\bullet$ we obtain a new factorization.
Since $(L')^\bullet$ agrees with $L^\bullet$ except in degree $c$
and since $(L')^c = \bigoplus_{\lambda \in \Lambda'} R$ the
induction step is proved.

\medskip\noindent
The conclusion of the discussion of the preceding paragraph is that
$1_K : K \to K$ factors as
$$
K \xrightarrow{\varphi} L \xrightarrow{\psi} K
$$
in $D(R)$ where $L$ can be represented by a finite
complex of free $R$-modules. In particular we see that $L$ is
perfect. Note that $e = \varphi \circ \psi \in \text{End}_{D(R)}(L)$
is an idempotent. By Derived Categories,
Lemma \ref{derived-lemma-projectors-have-images-triangulated}
we see that $L = \Ker(e) \oplus \Ker(1 - e)$.
The map $\varphi : K \to L$ induces an isomorphism with
$\Ker(1 - e)$ in $D(R)$. Hence we finally conclude that
$K$ is perfect by Lemma \ref{lemma-summands-perfect}.
\end{proof}

\end{document}
